\section{Introduction}

Let $\pi:X\to X_0$ and $\pi^!:X^!\to X_0^!$ be a symplectic-dual pair
of conical symplectic resolutions, equipped with Hamiltonian actions of
tori $T$ and $T^!$. Assume that the fixed-point sets are finite.
Symplectic duality supplies a bijection
\[
  X^T\xrightarrow{\sim}(X^!)^{T^!}.
\]
After choosing generic cocharacters $\sigma$ and $\eta$, stable
envelopes give isomorphisms
\[
  \HH_T(X)_\sigma
  \xleftarrow{\sim}W_{\sigma,\eta}
  \xrightarrow{\sim}
  \HH_{T^!}(X^!)_\eta.
\]
The two specialized equivariant cohomology groups are filtered by
cohomological degree. Transporting these filtrations to
$W_{\sigma,\eta}$ produces a bifiltration, whose bigraded Hilbert series
is the \emph{bicharacteristic polynomial} $T_{X,X^!}(x,y)$. This
construction was introduced by Davison and McBreen in their study of
symplectic duality and the Tutte polynomial \cite{mcbreen2025}; it is
recalled precisely in \Cref{sec:bichar_poly}.

In this paper we compute the bicharacteristic polynomial for the
self-dual type $A$ Springer resolution
\[
  \pi:T^*\Fln=\widetilde{\mathcal N}\longrightarrow\mathcal N.
\]
Its torus-fixed points are indexed by $S_n$, and its Steinberg variety
gives the convolution algebra
\[
  \HH^{BM}_{\mathrm{top}}
    (\widetilde{\mathcal N}\times_{\mathcal N}
     \widetilde{\mathcal N})
  \cong\mathbb C[S_n].
\]
Nilpotent orbits, Springer representations, and the terms in the
semismall decomposition are indexed by partitions of $n$. These two
descriptions meet in the following formula.

\begin{theorem}[thm:bichar_full_flag]{}
  For the self-dual type $A$ Springer resolution,
  \[
    T_{X,X^!}(x,y)
    =\sum_{\lambda\vdash n}
      K_{\lambda,(1^n)}(x)K_{\lambda^T,(1^n)}(y).
  \]
\end{theorem}

The principal issue in the proof is compatibility between the
stable-envelope isomorphism and the Springer decomposition. A splitting
provided by the decomposition theorem is not canonical, so we formulate
this compatibility using the support filtration. Category $\mathcal O$
identifies its increasing pieces with spans of projective classes and
its decreasing pieces with spans of simple classes. Koszul duality
exchanges these bases. Together with the fixed-point correspondence
$w\mapsto w^{-1}w_\circ$ and the Robinson--Schensted correspondence,
this shows that the stable-envelope isomorphism sends
\[
  F_{\leq\lambda}\longrightarrow F^!_{\geq\lambda^T}.
\]
It therefore induces, on the associated graded, an isomorphism
\[
  \IH^\bullet(\overline{\mathcal O_\lambda})
    \otimes\Htop(\mathcal B_\lambda)
  \longrightarrow
  \IH^\bullet(\overline{\mathcal O_{\lambda^T}})
    \otimes\Htop(\mathcal B_{\lambda^T}).
\]

The original and transported Springer actions commute: after
localization, one is right multiplication on the fixed-point set
$S_n$, while the other is left multiplication. A linear-algebra
argument then shows that the preceding map exchanges the two tensor
factors. Finally, Lusztig's formula for the intersection cohomology of
nilpotent-orbit closures identifies their graded dimensions with the
Kostka--Foulkes polynomials appearing in the theorem.

The formula immediately implies
\[
  T_{X,X^!}(x,y)=T_{X,X^!}(y,x)
  \qquad\text{and}\qquad
  T_{X,X^!}(1,1)=n!.
\]
The case $n=3$ is worked out explicitly at the end of
\Cref{chapter_main_thm}.

\paragraph{Organization.}
\Cref{chapter_combinatorics} fixes the combinatorial conventions.
\Cref{chapter_symplectic_duality} recalls stable envelopes and defines
the bicharacteristic polynomial. The Springer and Steinberg geometry is
collected in \Cref{chapter_springer}, and the semismall decomposition
and Kostka formula in \Cref{chapter_bbdg}. The proof of the theorem is
given in \Cref{chapter_main_thm}.
